\chapter{Entropy, and Why Determinism Is Not Enough}
\label{ch:ent}
% ===================================================================

Everything built so far is Turing complete and nothing built so far is
more than that.

It is worth being blunt about this, because the machinery has been
elaborate enough to disguise it.
The rho calculus is Turing complete.
Metering it does not add power --- the cost monad of
Chapter~\ref{ch:costmonad} takes a theory and returns a theory that does
strictly less, since a term which cannot pay does not fire.
Logging it does not add power either; the history monad records what
happened and computes nothing new.
Decorating rewrites with semiring values, imposing reversibility,
imposing arbitrage-freedom: all of these carve out subcategories, and a
subcategory of a category of Turing-complete theories contains only
Turing-complete theories.
Part~\ref{part:physics} descended a ladder and never once left the room.

And yet the learner of Part~\ref{part:ecology} does something that a
budget-constrained Turing machine does not obviously do.
It \emph{picks}.
When two of its sends contend for one receive, one of them wins, and
nothing in the term determines which.
This chapter is about that gap: where it is, what fills it, and why the
answer has the same shape as entropy.

\section{The closed room}

Recall the two constructions of Part~\ref{part:machinery} and what each
does to expressive power.

The cost monad installs a meter at the interaction cut.
Its effect on the transition system is subtractive: transitions that were
available become unavailable when the stack is exhausted.
Chapter~\ref{ch:costmonad} makes this precise as graded adequacy --- the
metered theory is adequate for the unmetered one up to the grading, which
is to say it distinguishes no more and sometimes less.

The history monad reifies rewrite events as terms.
Its effect is additive but not powerful: a term can now name a step it
took, which is exactly what Chapter~\ref{sec:reversibility} needs to
build the time-symmetric envelope, but naming a step one has taken is
not computing anything one could not compute before.

So neither construction moves a theory out of its Turing degree, and the
axioms of Part~\ref{part:physics} are restrictions rather than
extensions.
The ladder of that part is a ladder within a single room.

\begin{remark}[Why this is not a defect to be patched]
\label{rmk:ent-not-a-defect}
The temptation at this point is to bolt on an oracle: declare that some
distinguished term has access to a halting decision, and proceed.
That would be a definition, not a discovery, and it would leave entirely
open the question of what such a term is doing physically.
The route taken here is the opposite one.
We look for a place in the machinery that is \emph{already}
underdetermined --- a place where the theory as stated does not say what
happens next --- and ask what would have to be supplied there.
The claim of this part is that there is such a place, that it is not
exotic, and that a cost-accounted concurrent calculus is full of them.
\end{remark}

\section{Dissipation is where information leaves}

Chapter~\ref{sec:cost} established that a cost-accounted theory is
frictionful by design and therefore lives in the dissipative regime:
every lossy engine that fires strictly decreases the free energy, and the
virtual token is a Lyapunov function rather than a conserved charge.

Read that result backwards.
A frictionless theory is one whose engines are exactly value-preserving,
and by Theorem~\ref{thm:vtok-ftap} such a theory has an exact potential,
which is to say it has enough structure to say where every unit of value
came from.
A frictionful theory does not.
The dissipation $\theta$ accumulated around a cycle is precisely the
amount of provenance the theory cannot reconstruct: value went somewhere,
and the ledger does not record where.

That is an information-theoretic statement wearing economic clothing, and
Chapter~\ref{sec:cost} already turns it into one, by way of the Landauer
bound: the erasure a computation performs is bounded below by the
holonomy of the energy form.
Erasure is forgetting which of several possible predecessors a state
had.
Dissipation is the price of forgetting.
The two thresholds coincide, which is the most satisfying result in that
chapter and is also, for present purposes, a signpost.

\begin{remark}[Two erasures, and the gap between them]
\label{rmk:two-erasures}
There are in fact two distinct erasures in play, and they are not the
same map~\cite{TwoErasures}.
One forgets the \emph{history}: the history monad's counit, which
discards the log and keeps the state.
The other forgets the \emph{cost}: the cost monad's counit, which
discards the meter and keeps the computation.
Neither factors through the other, and the failure to factor is graded
--- there is a lattice of intermediate laxities between them, each
corresponding to a theory that remembers some of one and some of the
other.

The gap between the two erasures is not a technical annoyance.
It is a supply of positions.
A theory sitting strictly between them knows something about its own past
that it cannot express as a cost and something about its own costs that
it cannot express as a history, and a computation with information of
that shape is exactly a computation with something to decide.
\end{remark}

\section{Choice is what fills the gap}

The place where a concurrent theory does not say what happens next is the
race.
When two sends contend for a single receive, the reduction relation
offers both continuations and the calculus is silent about which occurs.

The traditional reading is that this is harmless underspecification, to be
resolved by an implementation and forgotten.
The reading taken here, and developed in Chapters~\ref{ch:apr}
and~\ref{ch:agy}, is that the resolution is a mathematical object with a
measurable strength.
Picking a winner from each contending family is a choice function on that
family.
Asserting that a complete schedule exists --- that every race, for the
whole run, gets resolved --- is an instance of the axiom of choice, and
which instance depends on the size and structure of the families
involved.
Finite contention is free.
Contention indexed by $\omega$ requires countable or dependent choice.
Class-sized families require the full axiom.

So the extra power is not bolted on.
In any calculus where interaction is primitive and confluence fails, it
is already present at every cut, wearing the name \emph{scheduler}.
Chapter~\ref{ch:apr} distinguishes the cuts at which this is observable
--- the \emph{apertures} --- from the cuts at which it is not, which is
the substantive part of the analysis and the reason the chapter is called
what it is.
Chapter~\ref{ch:chs} then makes the correspondence a typing discipline:
choice principles are the types, and schedulers are the terms.

\section{Entropy and choice as one dial}

We can now say what this chapter's title means.

Entropy production, in the sense of Chapter~\ref{sec:cost}, measures how
much which-way information the dynamics discards per step.
Choice strength, in the sense of Chapter~\ref{ch:apr}, measures how much
must be supplied to say which way it went.
These are the same quantity approached from opposite sides, and the
conjecture that they are related quantitatively is the natural one.

\begin{conjecture}[Dissipation bounds choice]
\label{conj:dissipation-choice}
Let $c$ be a cut in a cost-accounted interactive theory, and let
$\theta(c)$ be the dissipation charged to the interaction at $c$.
Then the Weihrauch degree of the scheduling problem at $c$ is bounded
below by a function of $\theta(c)$: a cut that dissipates nothing carries
no observable choice, and a cut whose scheduling problem requires
dependent choice must dissipate at least the corresponding erasure cost.
\end{conjecture}

The forward half of this is close to available.
A frictionless cut is one whose engine is value-preserving in both
directions, and a cut which is reversible in that sense is not an
aperture by Definition~\ref{ndo:def:aperture}: there is nothing to
observe about which way it went, because it can be run backwards.
The converse half is open, and the obstruction is that the Weihrauch
lattice is not linearly ordered while the dissipation is a single real
number, so no function of $\theta$ alone can pin down a degree.
The honest statement is probably that $\theta$ bounds the degree from
below along each chain, which is weaker and may be all that is true.
Chapter~\ref{ch:chs} grades the correspondence in the Weihrauch lattice
and is the place to look for what is actually established.

\section{What this buys, and what it does not}

It does not buy hypercomputation for free.
Nothing in this chapter exhibits a term that decides the halting problem,
and the tower of Chapter~\ref{sec:fibration} remains a construction whose
inhabitants are hypothesised rather than built.

What it buys is a location.
If a physical system is to exceed what a Turing machine can do, the
excess has to be spent somewhere, and this chapter says where: at the
apertures, in the resolution of races, paid for in dissipated free
energy.
That is a considerably more constrained claim than the usual ones, and it
is falsifiable in the direction that matters --- if it turned out that
every aperture in a physically realisable system carried only finite
contention, then the framework would predict that no physical system
exceeds the Turing computable, and the question would be settled the
other way.

There is also a consequence for Part~\ref{part:ecology} that will be made
quantitative in Part~\ref{part:origins}.
A population of learners has apertures wherever two of its members
contend, and the number of such places grows with the structural
elaboration of the population rather than with its raw size.
A large uniform population has few apertures; a small deeply nested one
has many.
Whatever measures nesting therefore measures the population's capacity to
be more than the sum of its computations --- and Part~\ref{part:origins}
argues that the thing which measures nesting is the assembly index.

% ===================================================================
